\section{Artifact Evaluation and Proof Engineering}

\subsection{Verified Artifact Facts}

The following facts were checked locally in the repository before drafting.
The command \texttt{lake build} completed successfully:

\begin{verbatim}
Build completed successfully (811 jobs).
\end{verbatim}

The final theorem and axiom footprint were checked with:

\begin{lstlisting}
import FriedbergMuchnik.Main

#check FriedbergMuchnik.friedberg_muchnik
#print axioms FriedbergMuchnik.friedberg_muchnik
\end{lstlisting}

Lean reported that \texttt{friedberg\_muchnik} depends on
\texttt{propext}, \texttt{Classical.choice}, and \texttt{Quot.sound}.  No
additional project axiom appears in the final theorem's footprint.

\begin{table}
\centering
\begin{tabular}{ll}
\toprule
Item & Value \\
\midrule
Lean toolchain & \texttt{leanprover/lean4:v4.26.0} \\
Mathlib revision & \texttt{2df2f0150c275ad53cb3c90f7c98ec15a56a1a67} \\
Lean files under \texttt{FriedbergMuchnik} plus root & 17 \\
Lean source lines, including comments/blank lines & 4024 \\
Top-level declarations counted by regex & 178 \\
\texttt{lake build} & succeeds locally \\
Final theorem axiom footprint & \texttt{propext}, \texttt{Classical.choice}, \texttt{Quot.sound} \\
\texttt{sorry}/\texttt{admit}/escape-hatch grep & no source declarations found \\
Repository license & no \texttt{LICENSE} file found \\
CI workflow & no \texttt{.github} workflow found \\
\bottomrule
\end{tabular}
\caption{Artifact facts verified from the repository.}
\end{table}

\subsection{Proof-Engineering Lessons}

The development illustrates several representation lessons.

First, c.e. approximations should be represented as finite positive
information, not as prefix-fixed Boolean strings.  The latter representation
would obscure the basic monotonicity of enumeration and risks proving the
wrong theorem.

Second, use should be recorded by the evaluator itself.  Proving after the
fact that every successful computation has some finite use would add a
continuity argument at exactly the point where the priority construction
needs executable finite data.

Third, finite and infinite semantics should be separated.  The construction
acts only on finite snapshots and bounded runs; the final theorem speaks
about total characteristic-function oracles.  The bridge lemmas in
\texttt{InfiniteEval.lean} and \texttt{Approximation.lean} connect these
layers without mixing them.

Fourth, finite injury becomes tractable when the informal phrase "injured
finitely often" is replaced by a bounded rank.  The rank \(0,1,2\) is small
enough for arithmetic proof but expressive enough to capture appointment
and action.

Finally, the c.e. proof benefits from a second implementation of the stage
function.  The mathematical state is convenient for the priority proof; the
tuple state is convenient for primitive-recursive closure.  The simulation
theorem keeps these two views synchronized.

\begin{figure}
\centering
\begin{tikzpicture}[
  cell/.style={draw, minimum width=7mm, minimum height=7mm},
  arr/.style={-Latex, thick}
]
\node at (-0.7,0) {$O$};
\foreach \i/\v in {0/1,1/0,2/1,3/0,4/?,5/?} {
  \node[cell] (a\i) at (\i*0.7,0) {\v};
}
\node at (-0.7,-0.9) {$O'$};
\foreach \i/\v in {0/1,1/0,2/1,3/0,4/1,5/0} {
  \node[cell] (b\i) at (\i*0.7,-0.9) {\v};
}
\draw[thick, dashed] (2.45,0.45) -- (2.45,-1.35);
\node[above] at (1.05,0.45) {agree below use \(u\)};
\node[above] at (3.5,0.45) {may differ};
\node[draw, rounded corners, fill=blue!5, align=center] (run) at (1.75,-2.1)
  {\texttt{run k O c x = halt d}\\\texttt{d.use = u}};
\node[draw, rounded corners, fill=blue!5, align=center] (run2) at (1.75,-3.1)
  {\texttt{run k' O' c x = halt d}};
\draw[arr] (run) -- (run2);
\end{tikzpicture}


\caption{Use preservation: oracles may differ above the recorded use, but
agreement below it preserves the same bounded computation.}
\end{figure}

\begin{figure}
\centering
\begin{tikzpicture}[
  box/.style={draw, minimum width=24mm, minimum height=9mm, align=center},
  arr/.style={-Latex, thick}
]
\node[box, fill=gray!10] (hi) {higher priority\\preserved};
\node[box, fill=blue!8, right=4mm of hi] (sel) {least requiring\\attention};
\node[box, fill=red!8, right=4mm of sel] (lo) {lower priority\\initialized};
\node[box, fill=green!8, below=8mm of sel] (new) {updated state};
\draw[arr] (hi) -- (new);
\draw[arr] (sel) -- node[right]{appoint or act} (new);
\draw[arr] (lo) -- (new);
\end{tikzpicture}


\caption{One priority stage preserves higher-priority records, acts for
the least requirement needing attention, and initializes lower-priority
records.}
\end{figure}

